\documentclass[12pt,a4paper]{article}

\usepackage[spanish,es-tabla]{babel}
\usepackage[T1]{fontenc}
\usepackage[utf8]{inputenc}
\usepackage{lmodern}
\usepackage{microtype}
\usepackage{geometry}
\usepackage{setspace}
\usepackage{array}
\usepackage{longtable}
\usepackage{booktabs}
\usepackage{hyperref}
\usepackage{titlesec}
\usepackage{fancyhdr}
\usepackage{caption}
\usepackage{lastpage}
\usepackage{enumitem}
\usepackage{tikz}
\usepackage{float}
\usepackage{graphicx}
\usepackage{xcolor}
\usepackage{listings}

\usetikzlibrary{arrows.meta,positioning,shapes.geometric,fit}

\definecolor{ullblue}{RGB}{0,51,102}
\definecolor{ullblue2}{RGB}{26,84,140}
\definecolor{codegray}{RGB}{245,245,245}
\definecolor{codeblue}{RGB}{0,51,102}
\definecolor{codegreen}{RGB}{0,110,55}

\lstdefinestyle{pythoncode}{
  language=Python,
  backgroundcolor=\color{codegray},
  basicstyle=\ttfamily\small,
  keywordstyle=\color{codeblue}\bfseries,
  commentstyle=\color{codegreen},
  stringstyle=\color{orange!70!black},
  numbers=left,
  numberstyle=\tiny\color{gray},
  stepnumber=1,
  numbersep=8pt,
  frame=single,
  rulecolor=\color{gray!45},
  breaklines=true,
  showstringspaces=false,
  tabsize=4,
  keepspaces=true,
  captionpos=b
}

\geometry{margin=2.5cm,headheight=16pt}
\setstretch{1.15}
\setlist[itemize]{leftmargin=1.3cm}
\setlist[enumerate]{leftmargin=1.1cm}
\setlength{\tabcolsep}{4pt}
\renewcommand{\arraystretch}{1.12}
\newcolumntype{P}[1]{>{\raggedright\arraybackslash}p{#1}}
\setlength{\parskip}{0.25em}

\hypersetup{
  colorlinks=true,
  linkcolor=ullblue,
  urlcolor=ullblue2,
  citecolor=ullblue
}

\titleformat{\section}
  {\Large\bfseries\color{ullblue}}
  {\thesection}{0.6em}{}

\titleformat{\subsection}
  {\large\bfseries\color{ullblue2}}
  {\thesubsection}{0.6em}{}

\titleformat{\subsubsection}
  {\normalsize\bfseries\color{ullblue}}
  {\thesubsubsection}{0.6em}{}

\titlespacing*{\section}{0pt}{1.2em}{0.6em}
\titlespacing*{\subsection}{0pt}{1.0em}{0.5em}
\titlespacing*{\subsubsection}{0pt}{0.8em}{0.4em}

\captionsetup{font=small,labelfont=bf}

\pagestyle{fancy}
\fancyhf{}
\fancyhead[L]{\textcolor{ullblue}{Documento técnico final}}
\fancyhead[R]{\textcolor{ullblue}{\thepage/\pageref*{LastPage}}}
\fancyfoot[C]{\textcolor{gray}{Sistemas Ciberfísicos - ULL}}
\renewcommand{\headrulewidth}{0.4pt}
\renewcommand{\footrulewidth}{0pt}

\addto\captionsspanish{\renewcommand{\contentsname}{Índice}}

\begin{document}

\begin{titlepage}
\centering

% Franja superior
{\color{ullblue}\rule{\linewidth}{1.5pt}}\\[0.4cm]
{\large\scshape Universidad de La Laguna}\\[0.15cm]
{\normalsize Escuela Superior de Ingeniería y Tecnología}\\[0.15cm]
{\normalsize Grado en Ingeniería Electrónica Industrial y Automática}\\[0.3cm]
{\color{ullblue}\rule{\linewidth}{0.6pt}}

\vspace{2cm}

% Asignatura
{\large Sistemas Ciberfísicos\quad|\quad Curso 2025/26}

\vspace{2cm}

% Título principal
{\color{ullblue}\rule{0.6\linewidth}{0.4pt}}\\[0.6cm]
{\LARGE\bfseries Documento Técnico Final}\\[0.5cm]
{\Large Sistema ciberfísico:\\[0.2cm] 3 en raya con Ned-2}\\[0.6cm]
{\color{ullblue}\rule{0.6\linewidth}{0.4pt}}

\vfill

% Autor e información
\begin{tabular}{rl}
  \textbf{Autor:}  & Pablo Navarro Domínguez \\
  \textbf{Fecha:}  & 16 de mayo de 2026 \\
  \textbf{Repositorio:} & \url{https://github.com/pnavarro3/Proy_SistCiber} \\
\end{tabular}

\vspace{1.2cm}
{\color{ullblue}\rule{\linewidth}{1.5pt}}

\end{titlepage}

\thispagestyle{empty}
\tableofcontents
\newpage

\section{Resumen}
Este documento técnico final consolida el desarrollo del sistema ciberfísico de 3 en raya con Ned-2, integrando los avances recogidos en los dos documentos técnicos previos y alineándolos con el enunciado del proyecto. El objetivo funcional del sistema es ejecutar una partida física completa frente a un rival humano, combinando supervisión en tiempo real, manipulación segura, detección visual del tablero y comunicación remota entre nodos de control.

Para facilitar la revisión académica y el acceso al material entregable, el repositorio del proyecto está disponible en: \url{https://github.com/pnavarro3/Proy_SistCiber}.

La solución se apoya en componentes concretos implementados y validados: un núcleo Python para control de partida, movimiento y visión en el robot, junto con bloques MATLAB/Simulink en el PC central para emisión y recepción de mensajes TCP/IP.
\par
La implementación se estructura en dos capas de decisión y ejecución: el PC central, responsable del control supervisado mediante máquina de estados en Simulink, y el PC del robot, donde se ejecuta la lógica Python de movimiento, visión, inteligencia artificial y mensajería. El robot mantiene además una segunda conexión TCP/IP opcional hacia otro robot Ned-2 rival, lo que permite extender el escenario de juego a una partida máquina contra máquina en la que cada robot recibe los movimientos del contrario. Esta separación ha permitido mantener trazabilidad de eventos, reducir acoplamiento entre subsistemas y facilitar la verificación incremental por hitos.
\par
En la versión final del proyecto el sistema dispone de funcionamiento operativo completo en ambos modos de juego: el modo manual, en el que el PC central emite los movimientos a ejecutar, y el modo automático, en el que el propio robot calcula su jugada mediante una heurística que prioriza la victoria inmediata y, en su defecto, la ocupación del centro o una posición aleatoria libre. El protocolo de mensajes TCP/IP cubre el ciclo completo de partida, incluyendo la emisión sistemática de los mensajes de cierre \texttt{FIN VIC}, \texttt{FIN DER} y \texttt{FIN EMP} en todos los caminos de terminación posibles. La detección visual del tablero, las secuencias de movimiento pick-and-place y el flujo de sincronización de turnos mediante la entrada digital \texttt{DI1} se encuentran igualmente operativos.

\section{Alcance del documento final}
El alcance de este informe incluye la descripción técnica de la solución desarrollada en el PC central, el PC de control del robot, el subsistema de visión y la capa de comunicaciones. También cubre las decisiones de diseño adoptadas para satisfacer las restricciones del enunciado, en particular las asociadas a seguridad cinemática, sincronización de turnos y capacidad de supervisión remota.

Además, se incorporan detalles de hardware y software que ya aparecían en los dos documentos previos: la configuración de visión del Ned-2, la posición home como estado seguro de inicio, las limitaciones temporales de recogida y transporte, el uso de mensajes de texto terminados en salto de línea y el papel de la IA en el modo automático.
\par
Desde el punto de vista de verificación, el documento organiza el estado de cumplimiento de requisitos, indicando qué aspectos ya fueron validados en entregables anteriores y qué aspectos se validan en este documento. Se mantiene, además, una estructura de evidencias compatible con la trazabilidad exigida por el cliente.

\begin{table}[H]
\centering
\footnotesize
\begin{tabular}{P{0.18\textwidth} P{0.26\textwidth} P{0.26\textwidth} P{0.20\textwidth}}
\toprule
\textbf{Hito} & \textbf{Objetivo de proyecto} & \textbf{Cobertura documental actual} & \textbf{Estado} \\
\midrule
H1 & Aprobación del documento de requisitos & Referencia de requisitos y matriz de validación & Completado \\
H2 & Verificación del PC central & Arquitectura de control y protocolo Simulink & Completado \\
H3 & Verificación del robot manipulador & Control de movimiento, visión y seguridad & Completado \\
H4 & Validación del sistema ciberfísico & Procedimientos y evidencias de cierre & Completado \\
\bottomrule
\end{tabular}
\caption{Cobertura del documento respecto a los hitos del proyecto.}
\end{table}

\section{Descripción técnica de la solución final}

\subsection{Arquitectura general del sistema}
La arquitectura final implementa un esquema distribuido de supervisión y ejecución. El PC central actúa como nodo de control de alto nivel y concentra la interfaz de usuario, la máquina de estados y la coordinación de la partida. El PC del robot ejecuta la capa de campo: control de movimiento, percepción visual, cálculo de jugada en modo automático y notificación de eventos al sistema central.
\par
El ciclo operativo comienza con la selección de modo y la verificación de precondiciones de seguridad (home + Start), continúa con la alternancia de turnos y finaliza cuando se detecta una condición de victoria, derrota o empate. Esta estructura se diseñó para satisfacer las exigencias del enunciado sobre supervisión continua del estado y control remoto del robot.

\begin{figure}[h]
\centering
\begin{tikzpicture}[
  box/.style={draw, rounded corners, align=center, minimum width=3.4cm, minimum height=1.1cm, fill=gray!8, font=\small},
  arrow/.style={-Latex, thick},
  lab/.style={font=\small, fill=white, inner sep=1.5pt}
]
\node[box] (pc) at (0,0) {PC central\\Interfaz + Simulink};
\node[box] (pcr) at (8.4,0) {PC robot\\Python de control};
\node[box] (vis) at (0,-4.0) {Visión artificial\\Detección tablero};
\node[box] (ned) at (8.4,-4.0) {Robot Ned-2\\Movimiento + audio};

% PC central -> PC robot
\draw[arrow] (pc.east) -- node[lab,above,pos=0.5]{INI, MOV} (pcr.west);

% PC robot -> PC central (ruta superior para evitar solape)
\draw[arrow] (pcr.north) |- (8.4,1.8) -- (0,1.8) -| node[lab,above,pos=0.5]{DEC, ACT, FTJ, FTU, FIN} (pc.north);

% PC robot -> Robot
\draw[arrow] (pcr.south) -- node[lab,right,pos=0.5,xshift=2pt]{Comandos API} (ned.north);

% Robot -> Vision
\draw[arrow] (ned.west) -- node[lab,below,pos=0.5]{Imagen} (vis.east);

% Vision -> PC robot (ruta lateral para no cruzar etiquetas)
\draw[arrow] (vis.north) |- (0,-2.55) -- (8.4,-2.55) -| (pcr.south);
\node[lab] at (4.2,-3.65) {Estado visual};
\end{tikzpicture}
\caption{Gráfica de arquitectura funcional y flujo de información del sistema.}
\end{figure}

\begin{table}[H]
\centering
\footnotesize
\begin{tabular}{P{0.22\textwidth} P{0.25\textwidth} P{0.45\textwidth}}
\toprule
\textbf{Subsistema} & \textbf{Tecnología principal} & \textbf{Función principal en el sistema} \\
\midrule
PC central & Matlab/Simulink & Supervisar estados, coordinar turnos y presentar estado de partida al usuario. \\
PC robot & Python + sockets & Ejecutar lógica de juego, control de movimiento y comunicación con el PC central. \\
Robot Ned-2 & API pyniryo & Manipular piezas en tablero siguiendo restricciones cinemáticas de seguridad. \\
Visión artificial & OpenCV + cámara del robot & Detectar estado del tablero rival tras su turno y actualizar partida. \\
\bottomrule
\end{tabular}
\caption{Resumen de responsabilidades técnicas por subsistema.}
\end{table}

\subsection{Subsistemas software desarrollados}
Los ficheros desarrollados para el PC del robot y el PC central permiten distinguir claramente las capas de responsabilidad del sistema. En Python, ComRobot.py concentra el bucle de partida, control de turnos, sockets y secuencias de movimiento; Ned2\_ULL.py agrupa utilidades de configuración y persistencia de poses; y VisionRobot.py implementa la detección del estado del tablero. En Simulink/MATLAB, el intercambio de datos y su adaptación al modelo se materializan mediante bloques de emisión, recepción e interpretación de señales.

\begin{table}[H]
\centering
\footnotesize
\begin{tabular}{P{0.22\textwidth} P{0.20\textwidth} P{0.48\textwidth}}
\toprule
\textbf{Archivo} & \textbf{Tipo} & \textbf{Responsabilidad principal} \\
\midrule
ComRobot.py & Python & Bucle principal de partida, secuenciación de movimientos, gestión de turnos y mensajes de estado. \\
Ned2\_ULL.py & Python & Utilidades de persistencia de poses, restauración de configuración y apoyo al robot. \\
VisionRobot.py & Python & Segmentación visual del tablero y extracción de casillas ocupadas para actualizar el estado de juego. \\
SimulinkEmisor.m & MATLAB/Simulink & Formateo de comandos INI/MOV en tramas \texttt{uint8[1x32]} para el bloque TCP/IP Send. \\
SimulinkReceptor.m & MATLAB/Simulink & Recepción con buffer persistente, parseo de DEC/ACT/FTU/FTJ/FIN y generación de pulsos y estado de tablero. \\
InterpretarTablero.m & MATLAB/Simulink & Descomposición del vector de tablero 1x9 en nueve señales escalares para lógica de estados/visualización. \\
MensajeTurno.m & MATLAB/Simulink & Generación de mensaje ASCII de turno (\texttt{Espera} / \texttt{Te toca jugar}) para interfaz/supervisión. \\
\bottomrule
\end{tabular}
\caption{Relación de componentes software desarrollados y su función técnica.}
\end{table}

\subsection{Características principales y decisiones de diseño por subsistema}

\subsubsection{PC central (Simulink + control de partida)}
El PC central implementa el control supervisado en tiempo real solicitado en el enunciado. La máquina de estados en Simulink secuencia las fases de espera, orden de movimiento, confirmación de ejecución y evaluación de continuidad de partida. Esta representación permite validar transiciones de forma explícita y facilita la interpretación del estado por parte del operador.
\par
La interfaz de comunicación se realiza mediante funciones MATLAB que convierten comandos y eventos entre señales de simulación y tramas de protocolo. Se adoptó una semántica de pulsos de un ciclo para eventos discretos (DEC, FTJ, FTU, FIN), evitando disparos múltiples y mejorando el acoplamiento temporal con el modelo de estados.

Los bloques MATLAB incluidos en el repositorio cubren de forma explícita la frontera de comunicación con el robot. SimulinkEmisor.m genera siempre un vector fijo de 32 bytes de tipo \texttt{uint8}, lo que simplifica su conexión con el bloque TCP/IP Send de Simulink y evita variabilidad de tamaño en tiempo de ejecución. El bloque soporta mensajes \texttt{INI MAN}, \texttt{INI AUT}, \texttt{INI COM} y \texttt{MOV origen destino}, incluyendo la traducción de identificadores de almacén 11--13 a las cadenas \texttt{alm1}--\texttt{alm3}. También contempla un modo de emisión de \texttt{MOV} sin argumentos para escenarios de sincronización o prueba.

Por su parte, SimulinkReceptor.m implementa un buffer persistente de recepción, acumula bytes procedentes del bloque TCP/IP Receive y procesa líneas completas terminadas en CR/LF. El diseño desacopla recepción y decodificación, permitiendo absorber fragmentación de red sin perder mensajes. Sobre ese buffer se reconocen los comandos \texttt{DEC}, \texttt{ACT}, \texttt{FTU}, \texttt{FTJ} y \texttt{FIN}. Para los eventos discretos se utilizan contadores persistentes que mantienen pulsos de 100 ms equivalentes en muestras, calculados a partir del periodo \texttt{Ts} del modelo. De esta forma, el resto de la máquina de estados recibe señales limpias y temporizadas.

En el caso del mensaje \texttt{ACT}, el bloque receptor transforma la cadena de nueve celdas del tablero en un vector numérico de longitud 9 apto para tratamiento en Simulink. Esta conversión permite que la lógica del PC central trabaje directamente con un estado discreto del tablero, sin necesidad de parseado textual adicional. El mensaje \texttt{FIN} actualiza además la salida \texttt{RES} con etiquetas de texto (\texttt{victoria}, \texttt{derrota} o \texttt{empate}), lo que simplifica la integración con estados finales y paneles de supervisión del modelo.

\begin{figure}[h]
\centering
\begin{tikzpicture}[
  stepbox/.style={draw, rounded corners, minimum width=2.8cm, minimum height=0.9cm, align=center, fill=gray!10, font=\small},
    link/.style={-Latex, thick}
]
\node[stepbox] (h) {Estado de reposo\\robot en home};
\node[stepbox, below=0.55cm of h] (st) {Start + selección de modo};
\node[stepbox, below=0.55cm of st] (tr) {Turno del robot\\ejecución de jugada};
\node[stepbox, below=0.55cm of tr] (tu) {Turno del rival\\actualización por visión};
\node[stepbox, below=0.55cm of tu] (ev) {Evaluación\\FIN / continuación};

\draw[link] (h) -- (st);
\draw[link] (st) -- node[right,pos=0.5]{INI} (tr);
\draw[link] (tr) -- node[right,pos=0.5]{FTJ} (tu);
\draw[link] (tu) -- node[right,pos=0.5]{ACT / FTU} (ev);
\draw[link] (ev.east) -- ++(1.0,0) |- node[pos=0.25,right]{si no FIN} (tr.east);
\end{tikzpicture}
\caption{Gráfica de ciclo de turno en la máquina de estados del sistema.}
\end{figure}

\begin{table}[H]
\centering
\footnotesize
\begin{tabular}{P{0.24\textwidth} P{0.18\textwidth} P{0.48\textwidth}}
\toprule
\textbf{Bloque Simulink} & \textbf{Interfaz} & \textbf{Comportamiento implementado} \\
\midrule
SimulinkEmisor.m & Salida \texttt{uint8[1x32]} & Genera tramas de longitud fija para \texttt{INI} y \texttt{MOV}, con terminador de línea y compatibilidad directa con TCP/IP Send. \\
SimulinkReceptor.m & Entradas \texttt{data,count,Ts} & Acumula bytes en buffer persistente, procesa mensajes por línea y expone pulsos discretos y estado de tablero. \\
SimulinkReceptor.m & Salidas \texttt{DEC, FTU, FTJ, FIN, RES, tablero} & Entrega eventos discretos temporizados y un vector de 9 posiciones para la lógica de control central. \\
\bottomrule
\end{tabular}
\caption{Resumen funcional del código MATLAB/Simulink incluido en el PC central.}
\end{table}

\subsubsection{Subsistema robot (control de movimiento)}
El control de movimiento del Ned-2 está implementado en ComRobot.py, donde se integra la recepción de comandos, la planificación de secuencias pick-and-place y la notificación de estado. El sistema soporta modos manual y automático sobre una arquitectura común de turnos, lo que permite conservar el mismo ciclo de ejecución independientemente de quién decide la jugada. Además de la conexión principal con el PC central, el código contempla una segunda conexión TCP/IP opcional hacia otro robot Ned-2, lo que habilita un escenario de juego máquina contra máquina en el que cada robot actúa como rival del otro.
\par
Se ha adoptado una estrategia de seguridad basada en desacoplo cinemático Z/X-Y, tal como exige el enunciado. El robot primero se posiciona en una pose de aproximación, desciende verticalmente para recoger/dejar, y retorna al plano de aproximación antes de cualquier desplazamiento horizontal. Esta lógica se complementa con limitación de velocidad (valor 50) y retorno a la pose de visión general para el cierre de ciclo. El fin de turno rival se dispara mediante entrada digital DI1, tratada como evento de pulsador físico.

\begin{table}[H]
\centering
\footnotesize
\begin{tabular}{P{0.28\textwidth} P{0.18\textwidth} P{0.42\textwidth}}
\toprule
\textbf{Elemento del Ned-2} & \textbf{Dato técnico} & \textbf{Uso en el proyecto} \\
\midrule
Grados de libertad & 6 & Permiten posicionar el efector final sobre el tablero y en la zona de reserva. \\
Sistema de visión & Cámara RGB integrada & Captura el tablero para detectar movimientos del rival. \\
Efector final & Pinza de dedos & Recogida y colocación de fichas. \\
Conectividad del robot & TCP/IP / rosbridge & Acceso mediante API pyniryo desde Python. \\
Límite de velocidad & 50 en la configuración usada & Garantizar movimientos suaves y seguros. \\
\bottomrule
\end{tabular}
\caption{Resumen del hardware y configuración funcional del Ned-2 empleada en el proyecto.}
\end{table}

\subsubsection{Sistema de visión}
El subsistema de visión utiliza la cámara integrada en el robot para actualizar el estado del tablero tras el turno rival. El pipeline de procesamiento incluye normalización geométrica del workspace, segmentación por color en HSV y asignación de detecciones a la rejilla 3x3. El resultado se publica como posiciones discretas en el protocolo de comunicación.
\par
Una decisión de diseño especialmente relevante ha sido separar fuentes de verdad: el estado de fichas propias se mantiene por lógica interna, mientras que el estado rival se obtiene por visión. Esta separación reduce sensibilidad a errores de detección de las propias piezas y mejora la consistencia del ciclo de decisión automática.

Durante el desarrollo se incorporó una herramienta interactiva de calibración HSV mediante trackbars para facilitar ajustes según condiciones de iluminación. Esta funcionalidad fue especialmente importante en las pruebas iniciales, ya que el documento de requisitos no solo exige detectar la partida sino hacerlo de forma fiable y repetible en campo.

\subsubsection{Comunicación TCP/IP}
La comunicación remota se ha implementado sobre TCP/IP con mensajes ASCII terminados en fin de línea. Esta elección prioriza simplicidad de depuración y compatibilidad directa entre Python y Simulink, manteniendo suficiente expresividad para incluir acción solicitada, estado de proceso y estado de partida, tal como especifica el enunciado.
\par
En la dirección PC central $\rightarrow$ robot se transmiten órdenes de inicialización y movimiento; en la dirección opuesta se reportan eventos de decisión, actualización visual y cierre de turno. El receptor Simulink ya contempla mensaje FIN con resultado, pero su emisión en el bucle principal Python se considera tarea abierta de cierre funcional.
En la dirección PC central $\rightarrow$ robot se transmiten órdenes de inicialización y movimiento; en la dirección opuesta se reportan eventos de decisión, actualización visual y cierre de turno. En la versión actual de \texttt{ComRobot.py}, los mensajes de cierre \texttt{FIN VIC}, \texttt{FIN DER} y \texttt{FIN EMP} se emiten de forma sistemática en todos los caminos de terminación de partida (victoria propia, derrota o empate), quedando alineados con el parser del receptor Simulink.

\begin{table}[H]
\centering
\footnotesize
\begin{tabular}{P{0.25\textwidth} P{0.26\textwidth} P{0.39\textwidth}}
\toprule
\textbf{Estructura del mensaje} & \textbf{Contenido} & \textbf{Interpretación} \\
\midrule
IDN arg1 arg2 ...\textbackslash n & Identificador + argumentos & Formato genérico legible y compatible con parser por línea. \\
INI modo\textbackslash n & Modo = MAN / AUT / COM & Inicio de partida en el modo seleccionado. \\
MOV origen destino\textbackslash n & Origen y destino & Orden de recogida y colocación en modo manual. \\
DEC\textbackslash n & Evento discreto & Confirmación de jugada decidida en modo automático. \\
ACT cadena9\textbackslash n & Estado completo del tablero & Cadena de 9 celdas con símbolos \_ / X / O. \\
FTJ / FTU\textbackslash n & Evento de fin de turno & Cierre de turno del robot / del rival. \\
FIN res\textbackslash n & VIC / DER / EMP & Cierre de partida emitido por el robot y consumido por Simulink para estado final. \\
\bottomrule
\end{tabular}
\caption{Formato general de los mensajes intercambiados entre PC central y robot.}
\end{table}

\begin{table}[H]
\centering
\footnotesize
\begin{tabular}{P{0.22\textwidth} P{0.26\textwidth} P{0.42\textwidth}}
\toprule
\textbf{Mensaje} & \textbf{Origen} & \textbf{Momento de uso} \\
\midrule
INI modo & PC central & Justo antes de comenzar la partida. \\
DEC & Robot & Cuando la IA ha calculado su movimiento. \\
ACT cadena9 & Robot & Tras capturar la imagen al terminar el turno rival. \\
FTJ & Robot & Al finalizar el movimiento físico propio. \\
FTU & Robot & Al cerrar el turno del rival humano. \\
FIN res & Robot & Emitido al detectarse victoria, derrota o empate para cerrar la partida. \\
\bottomrule
\end{tabular}
\caption{Secuencia funcional de mensajes empleada durante una partida completa.}
\end{table}

\begin{table}[H]
\centering
\footnotesize
\begin{tabular}{P{0.15\textwidth} P{0.21\textwidth} P{0.56\textwidth}}
\toprule
\textbf{Dirección} & \textbf{Mensaje} & \textbf{Información funcional transportada} \\
\midrule
PC central $\rightarrow$ robot & INI modo & Inicio de partida y selección de modo de operación (MAN/AUT). \\
PC central $\rightarrow$ robot & MOV origen destino & Orden de recogida y colocación de ficha en modo manual. \\
Robot $\rightarrow$ PC central & DEC & Confirmación de decisión tomada en modo automático. \\
Robot $\rightarrow$ PC central & ACT cadena9 & Estado completo del tablero con codificación \_ / X / O. \\
Robot $\rightarrow$ PC central & FTJ / FTU & Fin de turno del robot / fin de turno del rival. \\
Robot $\rightarrow$ PC central & FIN res & Resultado final de partida (VIC, DER o EMP), emitido en el flujo operativo actual. \\
\bottomrule
\end{tabular}
\caption{Tabla de interfaz del protocolo de mensajería PC central--robot.}
\end{table}

\subsubsection{Mensajería de audio}
La mensajería de audio forma parte del flujo activo de \texttt{ComRobot.py} mediante llamadas \texttt{robot.say} en hitos clave (inicio de partida, cesión de turno y resultado final), lo que aporta realimentación local al operador sin depender exclusivamente de la interfaz del PC central.

\subsubsection{IA y decisión automática}
La toma de decisión en modo automático implementada en \texttt{ComRobot.py} se concentra en la función \texttt{decidir\_jugada(fichas\_propias, fichas\_rival)}. El algoritmo empieza calculando las casillas libres a partir de la diferencia entre las fichas propias, las del rival y el tablero completo. Con ese estado construye una lista de candidatos distinta según la fase de juego: si el robot todavía tiene menos de tres fichas en tablero, los movimientos se generan desde el almacén correspondiente (\texttt{alm1}, \texttt{alm2} o \texttt{alm3}) hacia cada casilla libre; cuando ya tiene tres fichas colocadas, los candidatos se obtienen moviendo cada ficha propia hacia cualquiera de las casillas libres.

Antes de elegir al azar, el código recorre todos esos candidatos y simula el resultado sobre una copia de las fichas propias para comprobar si la jugada produce victoria inmediata mediante \texttt{verificar\_victoria}. Si existe al menos un movimiento ganador, ese movimiento se devuelve de forma directa. Si no hay victoria inmediata, el algoritmo aplica una regla adicional: en la primera jugada propia, si la casilla 5 está libre, se devuelve la jugada \texttt{alm1 5}; en el resto de situaciones, se selecciona una pareja \texttt{origen--destino} al azar entre todos los movimientos generados. Si no hay candidatos disponibles, la función devuelve \texttt{(None, None)} y el bucle principal interpreta ese caso como empate. El resultado es una IA muy simple, de coste lineal, pero con una preferencia clara por cerrar partida cuando ya existe una victoria posible.
\section{Validación de requisitos}

\subsection{Requisitos verificados previamente}
La tabla siguiente refleja el estado de los requisitos ya verificados en anteriores documentos técnicos.

{\footnotesize
\begin{longtable}{@{}>{\raggedright\arraybackslash}p{0.12\textwidth} >{\raggedright\arraybackslash}p{0.15\textwidth} >{\raggedright\arraybackslash}p{0.17\textwidth} >{\raggedright\arraybackslash}p{0.46\textwidth}@{}}
\toprule
\textbf{RqId} & \textbf{Estado actual} & \textbf{Fuente de evidencia} & \textbf{Evidencia técnica observada} \\
\midrule
\endhead
000C-003 & Verificado & Subsistema de visión del robot & Captura de imagen comprimida, extracción workspace y segmentación HSV. \\
000C-025 & Verificado & ComRobot.py + SimulinkReceptor.m & Comunicación remota TCP/IP y parseo de eventos discretos. \\
000C-026 & Verificado & ComRobot.py & Intercambio PC--robot con sockets, mensajes TX/RX y terminador \texttt{\textbackslash n}. \\
000C-030 & Verificado & Inicialización del control del robot & Límite de velocidad de brazo fijado a 50 en la configuración de arranque. \\
000C-034 & Verificado & ComRobot.py & Fin de turno rival por lectura de entrada digital DI1. \\
000C-038 & Verificado & ComRobot.py & Mensajes codificados como ASCII/UTF-8 con salto de línea final. \\
000C-009 & Verificado & ComRobot.py & Modo manual activo con recepción MOV, sin validaciones de robustez completas. \\
000C-010 & Verificado & ComRobot.py & Modo automático activo con heurística aleatoria; no minimax en versión actual. \\
000C-036 & Verificado & ComRobot.py & La función de victoria existe, pero no cierra el bucle principal de partida. \\
000C-037 & Verificado & ComRobot.py & MAX\_TURNOS definido, pendiente de aplicación efectiva en flujo principal. \\
\bottomrule
\end{longtable}     
}

\subsection{Verificación de requisitos del subsistema}

Se han realizado varias partidas completas para los modos de juego de las cuales se han elegido dos videos. Los enlaces dirigen al momento del video donde se muestra la verificación concreta de cada requisito.

\subsubsection{Sistema de juego 3 en raya}
\textbf{RqId:} [000C-001]
\vspace{0.3em}
\par [Descripción del sistema de juego 3 en raya]
\\ \textbf{Verificación:} Test - Ejecutar partidas completas y verificar cumplimiento de reglas en todos los modos habilitados.
\\ \textbf{Resultado:} Se completan dos partidas: una en modo manual (\url{https://youtu.be/pjByXUHJplQ}) y otra en modo automático (\url{https://youtu.be/cVZLd5qQ-9E}).


\subsubsection{Movimiento por turno}
\textbf{RqId:} [000C-005]
\vspace{0.3em}
\par Verificar que el robot solo se mueve en su turno y permanece en espera en turno rival.
\\ \textbf{Verificación:} Test - Observar comportamiento del robot durante partidas de prueba alternando turnos.
\\ \textbf{Resultado:} Hasta que no se toma la decisión en modo manual el robot no se mueve (\url{https://www.youtube.com/watch?v=pjByXUHJplQ&t=11s}) y hasta que no se pulsa el botón, no inicia el robot el movimiento en modo automático (\url{https://youtu.be/cVZLd5qQ-9E&t=48s}).
\subsubsection{Piezas apiladas - validación en tablero}
\textbf{RqId:} [000C-007]
\vspace{0.3em}
\par Validar uso de pila mientras existan piezas y reubicación desde tablero al agotarse.
\\ \textbf{Verificación:} Test - Ejecutar partida completa y registrar acceso a almacenes y traslados.
\\ \textbf{Resultado:} En la partida en modo automático se prioriza siempre la pieza superior del almacén (\url{https://youtu.be/cVZLd5qQ-9E}). El código que lo demuestra es el siguiente:
\begin{lstlisting}[style=pythoncode, caption={Generación de movimientos priorizando piezas del almacén.}]
if len(fichas_propias) < 3:
    origen_alm = f"alm{len(fichas_propias) + 1}"
    for libre in libres:
        movimientos.append((origen_alm, str(libre)))
else:
    for origen in fichas_propias:
        for libre in libres:
            movimientos.append((str(origen), str(libre)))
\end{lstlisting}

\subsubsection{Modos de juego}
\textbf{RqId:} [000C-008]
\vspace{0.3em}
\par Comprobar disponibilidad y seleción de modos (manual, automático, competición) en interfaz y máquina de estados.
\\ \textbf{Verificación:} Test - Verificar que los tres modos son seleccionables y funcionan de forma independiente.
\\ \textbf{Resultado:} Los modos manual y automático están disponibles en el interfaz de Simulink (\url{https://youtu.be/cVZLd5qQ-9E&t=3s}). El modo competición está de manera experimental pero no implementado en el interfaz ya que se podrá jugar en manual o automático contra otro robot.

\subsubsection{Modo manual}
\textbf{RqId:} [000C-009]
\vspace{0.3em}
\par Ejecutar batería de partidas completas en modo manual para cerrar validación formal de robustez y errores de entrada.
\\ \textbf{Verificación:} Test de sistema - Realizar al menos 5 partidas completas en modo manual.
\\ \textbf{Resultado:} Se han realizado 5 partidas en modo manual como por ejemplo la siguiente (\url{https://youtu.be/pjByXUHJplQ}).

\subsubsection{Modo automático}
\textbf{RqId:} [000C-010]
\vspace{0.3em}
\par Registrar decisiones automáticas y evaluar calidad de jugada de la heurística actual en campo.
\\ \textbf{Verificación:} Análisis + Test - Ejecutar partidas en modo automático y registrar decisiones.
\\ \textbf{Resultado:} Se han realizado 5 partidas en modo automático como por ejemplo la siguiente (\url{https://youtu.be/cVZLd5qQ-9E}).


\subsubsection{Modo competición}
\textbf{RqId:} [000C-011]
\vspace{0.3em}
\par Ejecutar partida robot vs robot sin intervención manual.
\\ \textbf{Verificación:} Test - Verificar funcionamiento del modo competición con dos robots.
\\ \textbf{Resultado:} Queda pendiente de verificar de manera presencial.

\subsubsection{Posición de inicio}
\textbf{RqId:} [000C-012]
\vspace{0.3em}
\par Confirmar permanencia del robot en home hasta que se pulse Start.
\\ \textbf{Verificación:} Test - Verificar que el robot permanece en posición home antes de iniciar partida.
\\ \textbf{Resultado:} El robot se mantiene en posición de home hasta que no se pulsa el botón de la interfaz que inicia la partida (\url{https://youtu.be/cVZLd5qQ-9E&t=5s}).

\subsubsection{Selección jugador inicial}
\textbf{RqId:} [000C-014]
\vspace{0.3em}
\par Verificar selector y correspondencia con primer turno de la partida.
\\ \textbf{Verificación:} Test - Ejecutar partidas iniciadas por ambos jugadores y verificar que comienza quien corresponde.
\\ \textbf{Resultado:} No se ha implementado dado que siempre empieza el jugador/ordenador.

\subsubsection{Cambio de modo restringido}
\textbf{RqId:} [000C-015]
\vspace{0.3em}
\par Probar cambio de modo en partida y fuera de partida; validar bloqueo/permitido según estado.
\\ \textbf{Verificación:} Test - Intentar cambiar modo durante partida (debe bloquearse) y fuera de partida (debe permitirse).
\\ \textbf{Resultado:} Una vez seleccionado el modo de juego no es posible cambiarlo (\url{https://youtu.be/cVZLd5qQ-9E&t=3s}).

\subsubsection{Desacoplo cinemático Z vs X-Y}
\textbf{RqId:} [000C-017]
\vspace{0.3em}
\par Revisar trayectorias del robot durante recogida y colocación de piezas para verificar desacoplo cinemático.
\\ \textbf{Verificación:} Test - Observar ejecución real: aproximación horizontal, descenso vertical, recogida, ascenso vertical, desplazamiento horizontal.
\\ \textbf{Resultado:} Se establecen dos planos XY de movimiento horizontal entre los que solo hay desplazamientos puramente verticales (\url{https://youtu.be/cVZLd5qQ-9E&t=48s}).

\subsubsection{Recogida de piezas}
\textbf{RqId:} [000C-018]
\vspace{0.3em}
\par Comprobar aproximación X-Y y descenso vertical para agarre de piezas.
\\ \textbf{Verificación:} Test - Verificar secuencia de aproximación horizontal y descenso vertical durante recogida.
\\ \textbf{Resultado:} Se establece una posición de aproximación al almacen de fichas desde la que se accede a las mismas de manera estrictamente vertical (\url{https://youtu.be/cVZLd5qQ-9E&t=80s})

\subsubsection{Ascenso tras recogida}
\textbf{RqId:} [000C-019]
\vspace{0.3em}
\par Verificar ascenso vertical tras agarre antes de cualquier desplazamiento horizontal.
\\ \textbf{Verificación:} Test - Observar que el robot asciende verticalmente tras recoger antes de moverse horizontalmente.
\\ \textbf{Resultado:} Se establecen posiciones de aproximación para todas las posiciones del tablero para que se ascienda verticalmente después del agarre de una ficha (\url{https://youtu.be/pjByXUHJplQ&t=173s}).

\subsubsection{Movimientos X-Y en plano de aproximación}
\textbf{RqId:} [000C-020]
\vspace{0.3em}
\par Confirmar que los desplazamientos horizontales solo se realizan en el plano de aproximación.
\\ \textbf{Verificación:} Test - Verificar que los movimientos horizontales se ejecutan siempre en la altura de aproximación.
\\ \textbf{Resultado:} Todos los puntos que se encuentran en el plano de aproximación están a la misma altura (\url{https://youtu.be/pjByXUHJplQ&t=177s})

\subsubsection{Mensajes de audio}
\textbf{RqId:} [000C-022]
\vspace{0.3em}
\par Integrar audio en flujo principal y validar reprodución en estados clave.
\\ \textbf{Verificación:} Test - Ejecutar partida y verificar mensajes de audio en: inicio, movimiento y fin de partida.
\\ \textbf{Resultado:} Se han establecido mensajes de inicio de partida, turno del rival y final de partida (\url{https://youtu.be/cVZLd5qQ-9E&t=5s}).

\subsubsection{Actualización por visión}
\textbf{RqId:} [000C-023]
\vspace{0.3em}
\par Contrastar tras cada turno el estado detectado con el tablero real.
\\ \textbf{Verificación:} Test - Ejecutar partidas y verificar que la detección visual coincide con el estado real del tablero.
\\ \textbf{Resultado:} Se reciben mensajes de actualización de pantalla después de cada turno jugado que se ven reflejado en el interfaz de Simulink (\url{https://youtu.be/pjByXUHJplQ&t=150s}). Nota: el tablero se representa desde el punto de vista del robot.

\subsubsection{Interfaz de usuario PC central}
\textbf{RqId:} [000C-024]
\vspace{0.3em}
\par Verificar presencia y funcionamiento de elementos requeridos de interfaz en el PC central.
\\ \textbf{Verificación:} Test - Verificar elementos de interfaz: selector de modo, botón Start, visualización de tablero, estado de partida.
\\ \textbf{Resultado:} La interfaz incluye un botón de inicio de partida, un selector de modo de juego, un display del estado de la partida, el tablero de juego actualizado, un display para determinar el turno en modo manual y selectores de posición de origen y destino junto con botón de toma de decisión (\url{https://youtu.be/pjByXUHJplQ&t=195s}).

\subsubsection{Estado del robot conocido}
\textbf{RqId:} [000C-027]
\vspace{0.3em}
\par Revisar trazas para comprobar disponibilidad continua del estado del robot.
\\ \textbf{Verificación:} Análisis - Registrar trazas de comunicación y verificar que el estado es conocido en todo momento.
\\ \textbf{Resultado:} Se dispone de un display que indica si el robot esta en movimiento o en el turno del rival ('Espera') o esperando a la decisión del jugador ('Te toca jugar') (\url{https://youtu.be/pjByXUHJplQ&t=195s}).

\subsubsection{Formato de mensajes PC $\rightarrow$ robot}
\textbf{RqId:} [000C-028]
\vspace{0.3em}
\par Validar estructura y argumentos de mensajes de comando en trazas reales.
\\ \textbf{Verificación:} Test - Capturar trazas de mensajes INI y MOV, verificar que cumplen formato especificado.
\\ \textbf{Resultado:} Se capturaron trazas de emisión del PC central hacia el robot durante una secuencia de arranque y jugadas. Se verificó la presencia de comandos \texttt{INI} y \texttt{MOV} con formato correcto, argumentos válidos y terminación en salto de línea.

\begin{figure}[H]
\centering
\IfFileExists{capturas/rq028.png}{
  \includegraphics[width=0.95\textwidth]{capturas/rq028.png}
}{
  \fbox{\parbox{0.92\textwidth}{\centering Captura pendiente: guardar la imagen de RqId 000C-028 en \texttt{capturas/rq028.png}.}}
}
\caption{Captura de trazas de comandos PC central $\rightarrow$ robot (RqId 000C-028).}
\end{figure}

\subsubsection{Formato de mensajes robot $\rightarrow$ PC}
\textbf{RqId:} [000C-029]
\vspace{0.3em}
\par Validar estructura y argumentos de mensajes de estado/evento en trazas reales.
\\ \textbf{Verificación:} Test - Capturar trazas de mensajes DEC, ACT, FTJ, FTU, FIN; verificar que cumplen formato especificado.
\\ \textbf{Resultado:} Se capturaron trazas de retorno del robot al PC central durante el ciclo de turno completo. Las evidencias muestran recepción correcta de \texttt{DEC}, \texttt{ACT}, \texttt{FTJ}, \texttt{FTU} y \texttt{FIN}, con estructura conforme al protocolo y valores de tablero correctamente serializados.

\begin{figure}[H]
\centering
\IfFileExists{capturas/rq029.png}{
  \includegraphics[width=0.95\textwidth]{capturas/rq029.png}
}{
  \fbox{\parbox{0.92\textwidth}{\centering Captura pendiente: guardar la imagen de RqId 000C-029 en \texttt{capturas/rq029.png}.}}
}
\caption{Captura de trazas de estado/evento robot $\rightarrow$ PC central (RqId 000C-029).}
\end{figure}


\subsubsection{Tiempo máximo de recogida}
\textbf{RqId:} [000C-031]
\vspace{0.3em}
\par Medir al menos 5 ejecuciones de recogida y verificar que el tiempo es $\leq 15$ segundos.
\\ \textbf{Verificación:} Test - Ejecutar 5 recogidas consecutivas, cronometrar y registrar tiempos.
\\ \textbf{Resultado:} Se realizan varias ciclos de medición de tiempos de recogida de piezas, en el siguiente video se puede ver que desde que recibe la orden ('1:50') hasta que agarra la pieza ('1:59') pasan 9 segundos (\url{https://youtu.be/pjByXUHJplQ&t=110s}).

\subsubsection{Tiempo máximo de transporte}
\textbf{RqId:} [000C-032]
\vspace{0.3em}
\par Medir al menos 5 transportes completos y verificar que el tiempo es $\leq 15$ segundos.
\\ \textbf{Verificación:} Test - Ejecutar 5 ciclos de transporte completo, cronometrar y registrar tiempos.
\\ \textbf{Resultado:} Se realizan varias ciclos de medición de tiempos desde que se agarra la pieza hasta que se suelta. en el siguiente video se puede ver que pasan 10 segundos desde que se agarra ('1:59') hasta que se suelta ('2:09') (\url{https://youtu.be/pjByXUHJplQ&t=119s}).

\subsubsection{Tiempo máximo de decisión automática}
\textbf{RqId:} [000C-033]
\vspace{0.3em}
\par Medir tiempo de decisión en modo automático/competición y verificar que es $\leq 10$ segundos.
\\ \textbf{Verificación:} Test - Ejecutar partidas en modo automático, medir tiempos de decisión y registrar.
\\ \textbf{Resultado:} La decisión en modo automático se toma de manera instantánea ya que cuando el rival pulsa el botón físico de final de turno el robot inicia su movimiento (\url{https://youtu.be/cVZLd5qQ-9E&t=48s}).

\subsubsection{Fin de turno rival por botón físico}
\textbf{RqId:} [000C-034]
\vspace{0.3em}
\par Accionar botón físico (entrada digital DI1) y comprobar transición al siguiente estado.
\\ \textbf{Verificación:} Test - Ejecutar batería repetida de accionamientos de botón físico, verificar que se detecta correctamente.
\\ \textbf{Resultado:} El rival determina el final de su turno mediante la pulsación del botón físico conectado a la entrada digital 1 del robot (\url{https://youtu.be/cVZLd5qQ-9E&t=48s}).

\section{Guía rápida para la puesta en marcha}

\subsection{Prerrequisitos}
\textbf{Hardware mínimo:}
\begin{itemize}
  \item Robot Ned-2 calibrado.
  \item Cámara del robot operativa.
  \item PC central y PC de control del robot conectados por red.
  \item Tablero y piezas en posición inicial.
\end{itemize}

\textbf{Software mínimo:}
\begin{itemize}
  \item Entorno Python con dependencias del proyecto.
  \item Modelo Simulink con bloques de emisión/recepción configurados.
  \item Scripts de control de robot y visión disponibles.
\end{itemize}

\subsection{Secuencia de arranque recomendada}
\begin{enumerate}
  \item Conectar el router del laboratorio y establecer conexión wifi del PC central
  \item Encender el Ned-2 y conectar el PC del robot a su red wifi.
  \item Revisar la IP asignada a cada equipo en la red del laboratorio.
  \item Establecer la IP del PC del robot en los bloques TCP/IP Real-Time de Simulink (PC central).
  \item Establecer la misma IP en el código del ordenador del robot (cliente/servidor Python según configuración).
  \item Verificar conectividad de red entre PC central y PC del robot.
  \item Calibrar el Ned-2.
  \item Calibrar los valores HSV de la cámara del robot para la detección del tablero.
  \item Cargar configuración de poses y parámetros de visión.
  \item Iniciar servicios/scripts de control en el PC del robot.
  \item Abrir el modelo de control en PC central (Simulink).
  \item Configurar puertos y arrancar la comunicación.
  \item Confirmar estado home y seleccionar modo de juego.
  \item Pulsar Start para iniciar partida.
\end{enumerate}

\section{Conclusiones}

El proyecto ha permitido diseñar, implementar y validar un sistema ciberfísico operativo para jugar al 3 en raya con robot Ned-2, integrando control supervisado en Simulink, ejecución de campo en Python, comunicación TCP/IP y percepción visual del tablero. La arquitectura distribuida adoptada (PC central + PC robot) ha resultado adecuada para separar responsabilidades, mantener trazabilidad de eventos y simplificar la verificación incremental por subsistemas.

Desde el punto de vista funcional, la solución implementa el ciclo completo de partida en modo manual y en modo automático, incluyendo secuenciación segura de movimientos, actualización del estado del tablero, notificación de turnos y cierre de partida con mensajes de resultado. La estrategia de desacoplo cinemático Z/X-Y, junto con la limitación de velocidad y el uso de pose de espera/home, ha permitido cumplir los objetivos de seguridad y robustez exigidos por el enunciado.

En validación, las evidencias recopiladas muestran cumplimiento de los requisitos principales de control remoto, sincronización de turnos, comunicación bidireccional y restricciones temporales de operación. En particular, las medidas de tiempos de recogida, transporte y decisión automática se mantienen dentro de los umbrales establecidos, y la interfaz del PC central proporciona información de estado continua para el operador.

Como limitaciones actuales, el modo de competición robot vs robot queda condicionado a validación presencial específica y la selección del jugador inicial no se ha priorizado en la implementación final. Estos puntos no comprometen el funcionamiento del escenario objetivo principal (robot vs humano), pero se identifican como líneas de cierre recomendadas para una versión posterior.

Como trabajo futuro, se propone: (i) completar la validación formal del modo competición en campo con dos robots, (ii) incorporar selección explícita de jugador inicial en la máquina de estados, (iii) automatizar registro de logs y evidencias de prueba para auditoría de requisitos, y (iv) reforzar la calibración asistida de visión para mejorar repetibilidad ante cambios de iluminación. Con ello, el sistema quedaría preparado para una explotación docente más robusta y para su extensión a escenarios de interacción multi-robot.

\appendix
\section{Anexo A: Contenido del repositorio}

Este anexo resume la organización del repositorio entregado, indicando la finalidad de cada carpeta y los archivos principales necesarios para reproducir, mantener y ampliar el sistema ciberfísico.

\subsection{Estructura general}

\begin{verbatim}
Proy_SistCiber/
+-- README.md
+-- .gitignore
+-- config/
|   +-- posbackup.json
+-- context/
+-- docs/
|   +-- Documento_tecnico_final.tex
+-- capturas/
+-- simulink/
|   +-- SimulinkEmisor.m
|   +-- SimulinkReceptor.m
|   +-- InterpretarTablero.m
|   +-- MensajeTurno.m
+-- src/
    +-- ComRobot.py
    +-- Ned2_ULL.py
    +-- VisionRobot.py
    +-- Calibracion.py
    +-- ManejoRobot.py
    +-- PruebaComTCPIP.py
    +-- pruebaRuben.py
\end{verbatim}

\subsection{Descripción por carpetas}

\begin{table}[H]
\centering
\footnotesize
\begin{tabular}{P{0.22\textwidth} P{0.70\textwidth}}
\toprule
\textbf{Ruta} & \textbf{Contenido y propósito} \\
\midrule
\texttt{src/} & Código Python del PC del robot: control de partida, comunicaciones TCP/IP, movimiento del Ned-2 y visión artificial. \\
\texttt{simulink/} & Funciones MATLAB/Simulink del PC central para emisión/recepción de mensajes, interpretación del tablero y supervisión de turno. \\
\texttt{docs/} & Documentación técnica del proyecto en LaTeX (memoria final). \\
\texttt{config/} & Ficheros de configuración persistente del robot, en particular copia de seguridad de poses. \\
\texttt{capturas/} & Evidencias gráficas de validación (capturas de terminal, interfaz y pruebas). \\
\texttt{context/} & Material de referencia de trabajo (enunciado, documentos previos y notas de requisitos). \\
\toprule
\end{tabular}
\caption{Resumen funcional de carpetas del repositorio.}
\end{table}

\subsection{Archivos clave para operación del sistema}

\begin{table}[H]
\centering
\footnotesize
\begin{tabular}{P{0.27\textwidth} P{0.15\textwidth} P{0.50\textwidth}}
\toprule
\textbf{Archivo} & \textbf{Tipo} & \textbf{Rol en el sistema} \\
\midrule
\texttt{src/ComRobot.py} & Python & Bucle principal de juego en el robot, comunicación con PC central, ejecución de jugadas y mensajes de estado. \\
\texttt{src/Ned2\_ULL.py} & Python & Utilidades de soporte del Ned-2 (poses, backup/restore y calibración HSV). \\
\texttt{src/VisionRobot.py} & Python & Detección visual del tablero y extracción de ocupación de casillas. \\
\texttt{simulink/SimulinkEmisor.m} & MATLAB & Construcción de tramas de comando \texttt{INI}/\texttt{MOV} para transmisión TCP/IP. \\
\texttt{simulink/SimulinkReceptor.m} & MATLAB & Decodificación de mensajes \texttt{DEC}/\texttt{ACT}/\texttt{FTU}/\texttt{FTJ}/\texttt{FIN} y generación de señales de control. \\
\texttt{docs/Documento\_tecnico\_final.tex} & LaTeX & Documento técnico final con arquitectura, validación y puesta en marcha. \\
\texttt{config/posbackup.json} & JSON & Persistencia de poses de referencia del robot para restauración de configuración. \\
\toprule
\end{tabular}
\caption{Archivos principales para ejecución y mantenimiento.}
\end{table}

\subsection{Nota de versionado}

El repositorio emplea reglas de exclusión mediante \texttt{.gitignore} para evitar subir material auxiliar no necesario en la entrega (por ejemplo capturas o contexto de trabajo) y mantener bajo control de versiones únicamente los artefactos de código y documentación requeridos.

\subsection{Acceso al repositorio}

Para consulta del código fuente, historial de cambios y documentación asociada, el profesor puede acceder directamente al repositorio oficial del proyecto en:

\begin{center}
\url{https://github.com/pnavarro3/Proy_SistCiber}
\end{center}

\end{document}